\let\cleardoublepage\clearpage
\part{Экономика, бизнес и услуги}
\let\cleardoublepage\clearpage
\chapter{Экономика, бизнес и услуги}
\ID{МРНТИ 06.52.13}{}

\begin{header}
\swa{Г.М.~Аубакирова, Ф.М.~Исатаева, С.К.~Мажитова, Л.И.~Тогайбаева, С.М.~Оспанова, А.Т.~Дюсекеева}{ЦЕНТРАЛЬНАЯ АЗИЯ В ГЛОБАЛЬНОЙ КОНКУРЕНЦИИ ЗА РЕДКОЗЕМЕЛЬНЫЕ МЕТАЛЛЫ: РИСКИ ЦЕПОЧЕК ПОСТАВОК И ПУТИ УКРЕПЛЕНИЯ ПОЗИЦИЙ}

\tsp{1}Г.М.~Аубакирова,
\tsp{1}Ф.М.~Исатаева\envelope,
\tsp{2}С.К.~Мажитова,
\tsp{1}Л.И.~Тогайбаева,
\tsp{3}С.М.~Оспанова,
\tsp{4}А.Т.~Дюсекеева
\end{header}

\begin{affil}
\tsp{1}Некоммерческое акционерное общество «Карагандинский технический университет Им. Абылкаса Сагинова», Казахстан,

\tsp{2}Карагандинский университет Казпотребсоюза, Караганда, Казахстан,

\tsp{3}ТОО «ONYX LS», Астана, Казахстан,

\tsp{4}Казахский университет технологии и бизнеса имени К. Кулажанова, Астана, Казахстан

\corrauthor{Корреспондент-автор: isataeva.farida@gmail.com}
\end{affil}

По мере ускорения глобального перехода к декарбонизации и цифровизации
растёт спрос на критически важные энергетические переходные минералы.
Прогнозируемый спрос на минералы и геоэкономическая фрагментация
становятся катализатором конкурентной борьбы за сохранение ресурсов,
привлекают новых игроков на инвестиционный рынок, меняют глобальные
цепочки поставок. В статье обосновано, что странам Центральной Азии,
обладающих значительными запасами критических минералов и расположенных
на перекрестке стратегических транзитных маршрутов, отведено особое
место в конкуренции крупнейших экономических держав вокруг отрасли
редкоземельных металлов. PESTEL-анализ позволил выявить ключевые угрозы
развития отрасли редкоземельных металлов, структурные ограничения и
внутренние уязвимости Центральной Азии, оценить возможности
центрально-азиатских стран для интеграции в высокотехнологичные сегменты
цепочек создания стоимости как внутри региона, так и за его пределами.
Исследованиями установлено, что для укрепления стратегической позиции
Центральной Азии среди системно значимых экономик с позиции управления
цепочками поставок критических минералов, необходимы структурные сдвиги
в пользу технологически сложных секторов. Авторы предлагают комплексный
подход
к~\href{https://www.sciencedirect.com/topics/earth-and-planetary-sciences/supply-chain-management}{управлению
цепочкой поставок}~минералов, включающий выстраивание сбалансированной
многовекторной внешней политики, привлечение качественных инвестиций,
региональную координацию и расширение внутрипроизводственных
перерабатывающих мощностей.

{\bfseries Ключевые слова}: критически важные минералы, редкоземельные
металлы, глобальные цепочки создания стоимости, Центральная Азия,
стратегия, перерабатывающие мощности.

\begin{header}
ОРТАЛЫҚ АЗИЯ СИРЕК ЖЕР МЕТАЛДАРЫ ҮШІН ӘЛЕМДІК БӘСЕКЕСТЕ: ЖЕТКІЗУ ТІЗБЕГІНІҢ ТӘУЕКЕЛДЕРІ ЖӘНЕ ПОЗИЦИЯЛАРДЫ НЫҒАЙТУ

\tsp{1}Г.М.~Әубәкірова,
\tsp{1}Ф.М.~Исабаева\envelope,
\tsp{2}С.К.~Мәжитова,
\tsp{1}Л.И.~Тоғайбаева
\tsp{3}С.М.~Оспанова,
\tsp{4}А.Т.~Дүйсекеева
\end{header}

\begin{affil}
\tsp{1}Әбілқас Сағынов атындағы Қарағанды техникалық университеті, Қарағанды, Қазақстан,

\tsp{2}Қазтұтынуодағы Қарағанды университеті, Қарағанды, Қазақстан,

\tsp{3}"ONYX LS" ЖШС, Астана, Қазақстан,

\tsp{4}К.Қулажанов атындағы Қазақ технология және бизнес университеті, Астана, Қазақстан,

e-mail: isataeva.farida@gmail.com
\end{affil}

Әлемдік деңгейде декарбонизация мен цифрландыруға көшу жеделдеген сайын,
энергетикалық маңызды өтпелі минералдарға сұраныс артып келеді. Пайдалы
қазбаларға болжамды сұраныс және геоэкономикалық бөлшектену ресурстарды
үнемдеу үшін бәсекелестікті күшейтіп, инвестициялық нарыққа жаңа
ойыншыларды тартып, жаһандық жеткізу тізбегін қайта құруда. Бұл мақалада
маңызды минералдардың айтарлықтай қорына ие және стратегиялық транзиттік
жолдардың қиылысында орналасқан Орталық Азия елдері сирек кездесетін жер
металдар өнеркәсібінің айналасындағы ірі экономикалар арасындағы
бәсекелестікте ерекше орын алатыны айтылады. PESTEL талдауы сирек
кездесетін жер металдар өнеркәсібінің дамуына төнетін негізгі
қауіптерді, құрылымдық шектеулерді және Орталық Азиядағы ішкі
осалдықтарды анықтады және Орталық Азия елдерінің аймақ ішінде де, одан
тыс жерлерде де құндылық тізбектерінің жоғары технологиялық
сегменттеріне интеграциялану әлеуетін бағалады. Зерттеу Орталық Азияның
маңызды минералдарды жеткізу тізбегін басқару тұрғысынан жүйелік маңызды
экономикалар арасындағы стратегиялық орнын нығайту технологиялық
тұрғыдан күрделі секторларға құрылымдық ауысуды қажет ететінін анықтады.
Авторлар минералды жеткізу тізбегін басқаруға кешенді тәсілді ұсынады,
оның ішінде теңгерімді, көпвекторлы сыртқы саясатты әзірлеу, жоғары
сапалы инвестицияларды тарту, аймақтық үйлестіру және ішкі өңдеу қуатын
кеңейту бар.

{\bfseries Түйін сөздер:} маңызды минералдар, сирек жер металдары, жаһандық
құн тізбегі, Орталық Азия, стратегия, қайта өңдеу қуаты.

\begin{header}
CENTRAL ASIA IN GLOBAL COMPETITION FOR RARE-EARTH METALS: SUPPLY CHAIN RISKS AND WAYS TO STRENGTHEN POSITIONS

\tsp{1}G.M.~Aubakirova,
\tsp{1}F.M.~Isataeva\envelope,
\tsp{2}S.K.~Mazhitova,
\tsp{1}L.I.~Togaybaeva,
\tsp{3}S.M.~Ospanova,
\tsp{4}A.T.~Dyusekeyeva
\end{header}

\begin{affil}
\tsp{1}Abylkas Saginov Karaganda Technical University, Karaganda, Kazakhstan,

\tsp{2}Karaganda University of Kazpotrebsoyuz, Karaganda, Kazakhstan,

\tsp{3}"ONYX LS" LLP, Astana, Kazakhstan,

\tsp{4}K. Kulazhanov Kazakh University of Technology and Business, Astana, Kazakhstan,

e-mail: isataeva.farida@gmail.com
\end{affil}

As the global transition toward decarbonization and digitalization
accelerates, demand for critical energy transition minerals is growing.
Projected demand for minerals and geoeconomic fragmentation are
catalyzing competition for resource conservation, attracting new players
to the investment market, and reshaping global supply chains. This
article argues that Central Asian countries, possessing significant
reserves of critical minerals and located at the crossroads of strategic
transit routes, hold a special place in the competition among major
economies around the rare earth metals industry. A PESTEL analysis
identified key threats to the development of the rare earth metals
industry, structural constraints, and internal vulnerabilities in
Central Asia, and assessed the potential of Central Asian countries to
integrate into high-tech segments of value chains both within the region
and beyond. The research found that strengthening Central
Asia' s strategic position among systemically important
economies in terms of managing critical minerals supply chains requires
structural shifts toward technologically sophisticated sectors. The
authors propose a comprehensive approach to mineral supply chain
management, including developing a balanced, multi-vector foreign
policy, attracting high-quality investment, regional coordina\-tion, and
expanding in-house processing capacity.

{\bfseries Keywords:} critical minerals, rare earth metals, global value
chains, Central Asia, strategy, processing capacity.

\begin{multicols}{2}
{\bfseries Введение.} В условиях растущего мирового спроса на критически
важные энергетические переходные минералы (CETM), в том числе на
редкоземельные металлы (РЗМ), обеспечение устойчивых и
конкурентоспособных по стоимости цепочек их поставок становится наиболее
востребованной стратегической задачей {[}1{]}.

Современный геополитический и геоэкономический ландшафт внес коррективы
в глобальные цепочки создания стоимости (ГЦСС): наблюдается их
реконфигурация, возрастают наукоемкость, степень локализации и
регионализации поставок промежуточного продукта, усугубляется влияние
глобальных санкций на страны-участники, меняется характер услуг,
появляются бизнес-функции, от которых непосредственно зависят
потребительская ценность товара и решение профильных задачи конкретного
производства {[}2{]}.

Для отслеживания вклада услуг и различных способов поставки в
добавленную стоимость разных отраслей, вводятся новые показатели
торговли добавленной стоимостью (TiVA), позволяющие объективно оценить
результативность и адаптивность динамичных цепочек поставок, в том числе
горизонтально распределенных сетевых структур, отражающих вклад
независимых участников ГЦСС в процесс создания стоимости конечного
продукта {[}3{]}.

В поиске ответов на вопросы касательно того, как и какой ценой будет
расширяться предложение и регулирование доступности CETM, определяющую
роль играет государство. Именно оно выступает посредником и поставщиком
ресурсов, содействует созданию актуальной инфраструктуры {[}4,5{]}.

Для многих стран высококонцентрированные цепочки поставок обостряют
геополитические конфликты, влияют на смежные перерабатывающие отрасли и
энергетический баланс. Эти вызовы обуславливают необходимость реализации
мер государственной политики, направленных на стимулирование эффективной
интеграции стран в ГЦСС, повышение их экономического суверенитета.

Особенно решающая роль государственной поддержки глобальных цепочек
поставок проявляется в развивающихся странах, которые стремятся сделать
из ключевых минералов надежный стратегический актив, ускоряющий их
экономический рост, расширяющий рамки сотрудничества, основанные на~
прозрачности, подотчетности и устойчивости.~

Осознавая риски для национальной безопасности, связанные с подавляющем
доминированием Китая, на долю которого приходится более 70\% мировой
добычи РЗМ, 90\% их разделения и переработки, и вводимых им экспортных
ограничений, мировые лидеры предпринимают шаги по увеличению внутреннего
производства и диверсификации цепочек поставок.

\texorpdfstring{ Одновременно внимание мировых держав
переключилось на один из самых быстрорастущих субрегионов в мире -
Центральную Азию. Объясняется это не только обострением конкуренции за
ключевые полезные ископаемые, но и перестройкой евразийской торговли,
возобновлением интереса к альтернативным транспортным маршрутам,
связывающим Восточную Азию, Европу и Ближний Восток, позиционированием
Центральной Азии как целостного экономического пространства, способного
привлекать качественные инвестиции.
}{ Одновременно внимание мировых держав переключилось на один из самых быстрорастущих субрегионов в мире - Центральную Азию. Объясняется это не только обострением конкуренции за ключевые полезные ископаемые, но и перестройкой евразийской торговли, возобновлением интереса к альтернативным транспортным маршрутам, связывающим Восточную Азию, Европу и Ближний Восток, позиционированием Центральной Азии как целостного экономического пространства, способного привлекать качественные инвестиции. }

Минеральный профиль которых сформирован геологическим картированием еще
советского периода (рисунок 1).
\end{multicols}

\fig{e/image2}[Рис.1 - Доля мировых запасов минеральных ресурсов стран
Центральной Азии]

\begin{multicols}{2}
Эти государства интегрированы в ГЦСС главным образом на этапе добычи и
первичной обработки, обладают разнообразной минерально-сырьевой базой и
существенным объемом мировых запасов (марганцевая руда- 38,6\%, хром
-30,07\%, свинец -20\%, цинк -12,6\%, титан -8,7\%, алюминий/бокситы
-5,8\%, медь -5,3\%, кобальт -5,3\%, молибден -5,2\%) на этапе добычи и
начальной обработки, рассматриваются преимущественно
как~\href{https://t.me/tengenomika/8850}{поставщики первичного сырья}~с
критически низкой долей технологической переработки {[}6{]}.

Осознавая стратегическую ценность контроля цепочек поставок CETM, эти
страны стремятся освоить новые стратегии добычи, переработки,
производства и конечного их использования, продвигать глубокую
переработку ценнейшего сырья на своей территории, а не просто развивать
модель «сырье в обмен на технологии» в сотрудничестве с развитыми
экономиками. Благодаря возросшему спросу на РЗМ, страны Центральной Азии
впервые могут превратить свою ресурсную зависимость в стратегический
актив, перейти от роли пассивного поставщика сырья к субъекту
международной политики.

Нарастающая конкуренция крупнейших держав за критические минералы,
открывает для стран Центральной Азии возможность использовать
геоэкономическое давление как рычаг для качественной промышленной
модернизации национальных экономик, импорта высоких технологий и
внедрения экологически ответственных методов добычи и стандартов ESG.

Таким образом, анализ трансформации ГЦСС в секторе критических минералов
позволяет констатировать переход от модели свободной рыночной
конкуренции к модели геоэкономической фрагментации. Установлено, что
ключевым риском для устойчивости поставок является сверхвысокая
концентрация перерабатывающих мощностей (более 90\% по ряду РЗМ) в одной
юрисдикции. Это создает объективные предпосылки для позиционирования
Центральной Азии как альтернативного узла ГЦСС, способного обеспечить
диверсификацию поставок в условиях глобального энергоперехода.

В свете сказанного целью исследования является оценка экономического
потенциала индустрии критически важных минералов в странах Центральной
Азии в контексте геостратегических перестроек цепочек их поставок,
выявление ключевых рисков для мировых цепочек поставок РЗМ, разработка
путей укрепления позиций этих стран в условиях многополярной
конкуренции.

Объект исследования - страны Центральной Азии, обладающие огромными
запасами критических минералов, рассматриваемые в контексте глобальных
цепочек поставок.

Предмет исследования - геополитическая и геоэкономическая конкуренция
ведущих стран мира, главным образом, США, ЕС и Китая, за доступ к РЗМ
Центральной Азии, связанные с этим риски уязвимости цепочек поставок и
возможности региона использовать глобальную конкуренцию для
стимулирования внутристрановой добавленной стоимости.

Гипотеза строится на том, что интеграция Центральной Азии в ГЦСС
обусловлена богатством и разнообразием минеральных ресурсов, но
затруднена~геополитической конкуренцией, неблагоприятными
институциональными рамками, недостатком инвестиций и слабыми позициями в
средней и конечной стадиях переработки.

Научная значимость. Исследование вносит вклад в осмысление
нового геополитического фронта - конкуренции за критические минералы в
эпоху энергетического перехода и деглобализации цепочек поставок.
Исследование углубляет понимание того, как Центральная Азия будучи
периферийным регионом, не имеющего выхода к морю, пытается совместить
ресурсный национализм с устойчивым развитием и технологической
модернизацией, повысить внутреннюю переработку и экспортную ценность
минеральных ресурсов.

Предложенные в статье пути укрепления позиций региона имеют практическую
ценность для государственной политики стран Центральной Азии, опыт
которых окажется полезным развивающимся странам в их стремлении
соответствовать инициативам энергетического перехода и реализации
проактивного подхода к~возобновляемым источникам энергии.

{\bfseries Материалы и методы}. Исследование выполнено в рамках смешанных
методов с преобладанием качественного анализа. Применены общие законы
формализации, логики и сравнения, использован смешанный подход,
сочетающий качественный и количественный анализ данных открытого доступа
о горнорудном секторе стран Центральной Азии, для интеграции результатов
применялся PESTEL - анализ.

Основные методы: сбор и критический анализ актуальных открытых данных по
запасам, добыче, экспорту и переработке РЗМ с использованием
национальных статистических комитетов стран ЦА за исследуемый период;
оценка цепочек поставок с акцентом на ключевые узкие места (концентрация
переработки в Китае, технологическая зависимость, экологические
ограничения); сравнительный анализ стратегий центрально-азиатских стран
касательно их продвижения в глобальных в цепочках поставок ключевых
минералов, текущих инициатив и мировых соглашений, возможностей
оптимизации цепочек поставок и внутрирегиональной переработки в
обозримом будущем.

Авторы признают ограничения данного исследования - недостаточная
транспарентность данных по запасам и переработке РЗМ в Центральной Азии,
геополитическая ангажированность части источников, высокая волатильность
геополитической обстановки, влияющая на инвестиционные решения. В
открытых источниках конкретные количественные данные о ключевых РЗМ по
изучаемым странам отсутствуют ввиду конфиденциального характера
информации. По некоторым РМЗ дается только качественная оценка, многие
доступные данные до сих пор носят фрагментированный характер. Такая
ситуация объясняется как историческим прошлым Центральной Азии (наследие
секретности), так и нарастающим геополитическими рисками, желанием
некоторых внешних глобальных игроков не допускать утечки ценной
информации.

Исходя из вышесказанного, во избежание проблем этического и юридического
характера, методология исследования базируется на сочетании
качественного и количественного анализа, дополненного более точными
исходными данными, приведенными в открытых источниках. Применение
указанного методологического подхода обусловлено тем, что сведения о
ежегодном обновлении запасов поступают, во-первых, крайне редко,
во-вторых, они чрезмерно зависимы от результатов разведки, что может
отразиться на достоверности поступающих данных.

Несмотря на вышесказанное, авторы надеются, что с помощью данного
исследования, имеющего аналитический и прикладной характер, будет
инициирован дискурс по изучению стратегических ресурсов для стран с
развивающейся экономикой, формирующих собственные бизнес-модели
укрепления позиций в ГЦСС.

{\bfseries Результаты обсуждения.} В условиях уменьшения географической
протяженности ГЦСС и концентрации ключевых бизнес-функций в промышленно
развитых странах, государства Центральной Азии пытаются найти баланс
между экономической устойчивостью и открытостью, продвинуться вверх по
цепочке создания стоимости, заложить основу для ГЦСС следующего
поколения {[}7{]}.

Чтобы занять позицию нового центра добычи полезных ископаемых, крупного
мирового поставщика РЗМ и редкоземельных элементов (РЗЭ), этим странам
необходимо осознать международную конкурентоспособность и глобальную
рыночную стоимость отрасли РМЗ, оценить свой промышленный потенциал,
смоделировать региональную интеграцию (таблица 1).

Анализ данных, представленных в таблице 1, свидетельствует о
значительном и диверсифицированном минерально-сырьевом потенциале стран
Центральной Азии в сфере РЗМ и критически важных минералов. Казахстан и
Кыргызстан обладают наиболее крупными подтверждёнными запасами РЗЭ,
Узбекистан и Туркменистан - существенными ресурсами лития и
сопутствующих элементов, Таджикистан - перспективными проявлениями
ниобия, тантала и сурьмы. Выявленные в 2024-2025 гг. новые участки и
упрощённые процедуры лицензирования указывают на активизацию
геологоразведочных работ и рост инвестиционной привлекательности
региона.
\end{multicols}

\tcap{Таблица 1- Минерально-сырьевая база критических металлов и РЗМ стран Центральной Азии}
\begin{longtblr}[
  label = none,
  entry = none,
]{
  width = \linewidth,
  colspec = {Q[154]Q[265]Q[623]},
  cells = {font = \small},
  row{1} = {c},
  cell{1}{1} = {font=\bfseries\small},
  cell{1}{3} = {font=\bfseries\small},
  cell{2}{1} = {c},
  cell{3}{1} = {c},
  cell{4}{1} = {c},
  cell{5}{1} = {c},
  cell{6}{1} = {c},
  hlines,
  vlines,
}
Страна       & {\textbf{Ключевые активы}\\\textbf{(месторождения)}}                                                                                 & Подтвержденные ресурсы / Типы РЗМ                                                                                                                                                                                                                                                                      \\
Казахстан    & Акбулак (РМЗ), Куйректыколь.                                                                                                         & Запасы РЗЭ оцениваются в 2,6 млн. тонн. Подтвержденные запасы месторождения Жана Казахстан (порядка 935 тыс. тонн). Выявлено 38 новых перспективных участков в 2024 г. В 2025 г. выставлено на продажу 40 месторождений, выдано более 3,2 тыс. лицензий, внедряется упрощённая система лицензирования. \\
Узбекистан   & Тебинбулак                                                                                                                           & Ресурсы лития (месторождение Нурликон 178,5 тыс. тонн). Значительные ресурсы вольфрама (потенциал 3-го места в мире по добыче). Ресурсы графита (Тасказган - никель-графитовое месторождение, свыше 1,3 млн тонн) [8,9]                                                                                \\
Кыргызстан   & Кутессай-II, Кызыл-Омполь                                                                                                            & {Иттрий, лантан, церий, титаномагнетит\\Ресурсы РЗМ (Кутессай-II, до 60 тыс. тонн). В 2024 г. снят мораторий на добычу урана и тория для стимуляции сопутствующей добычи РЗМ [10, 11]}                                                                                                                 \\
Таджикистан  & {60 проявлений РЗМ\\в Раштской долине (ниобий и тантал), 15 новых место-рождений~РМ, РЗЭ в районах Карасу, Агбасой, Пайрон, Рохшиф.} & Высокая концентрация сурьмы и лития. Наличие проявлений ниобия и тантала. В 2024 г. зафиксирован рост лицензирования участков РЗМ [12,13]                                                                                                                                                              \\
Туркменистан & Кара-Богаз-Гол (литий), легкие РЗЭ (неодим, лантан) в Лебапском велаяте.                                                             & Литий (рапа), йод, бром [14]                                                                                                                                                                                                                                                                           
\end{longtblr}

\begin{multicols}{2}
Такое распределение ресурсной базы создаёт объективные предпосылки для
формирования внутрирегиональной конфигурации цепочки создания стоимости,
в которой добыча сырья сосредоточена преимущественно в Таджикистане,
Кыргызстане и Казахстане, а стадии обогащения и первичной переработки -
в Казахстане и Узбекистане, обладающих относительно развитой
инфраструктурой и промышленными хабами. Данная конфигурация позволяет
минимизировать зависимость от внешних поставок на ранних стадиях ГЦСС и
повысить добавленную стоимость на последующих этапах (таблица 2).

В соответствии с методикой сравнительного анализа запасов РЗМ и РЗЭ в
странах Центральной Азии, авторы приходят к выводу, что среди государств
Центральной Азии по потенциалу РЗМ и РЗЭ наилучшие позиции занимают
Казахстан и Узбекистан, лидирующие также в инновационном развитии: в
2021-2025 гг. Казахстан занимал в мировом рейтинге 77-83 места, в 2025
г. Узбекистан опередил Казахстан, поднявшись с 93 на 79 позицию, также в
Узбекистане в 2025г. прогнозируется более высокий экономический рост
(рисунок 2) {[}15{]}.
\end{multicols}

\tcap{Таблица 2 - Конфигурация внутрирегиональной цепочки создания стоимости}
\begin{longtblr}[
  label = none,
  entry = none,
]{
  width = \linewidth,
  colspec = {Q[140]Q[602]Q[198]},
  cells = {font = \small},
  row{1} = {c,font=\bfseries\small},
  cell{2}{1} = {c},
  cell{3}{1} = {c},
  cell{4}{1} = {c},
  hlines,
  vlines,
}
Этап ГЦСС                & Роль стран-участниц                                                                                                            & Ожидаемая добавленная стоимость \\
Добыча                   & Таджикистан, Кыргызстан, Казахстан                                                                                             & Минимальная (сырьевой экспорт)  \\
Обогащение / Разделение    & Казахстан, Узбекистан (хабы по добыче и переработке полезных ископаемых в Навоийской, Ташкентской и Сурхандарьинской областях) & +40-60\% (концентраты, оксиды)  \\
Производство компонентов & Региональная кооперация, инвесторы (ЕС, Корея)                                                                                 & +150-300\% (магниты, сплавы)    
\end{longtblr}

\fig[0.9\textwidth]{e/image3}[Рис.2 - Реальный рост ВВП стран Центральной Азии, \%\\\normalfont{\emph{Источник: International Monetary Fund.2025. World Economic
Outlook:Global Economy in Flux, Prospects Remain Dim. Washington, DC.
October}}]

\begin{multicols}{2}
Новые разведанные месторождения закрепляют статус Казахстана, как
источника критически важных материалов, в целом страну можно
рассматривать как нетто-экспортера РЗМ (рисунок 3).

Несмотря на ориентацию инфраструктуры Казахстана в большей мере на
цветные металлы, а не на переработку РЗМ, страна продвигается в создании
перерабатывающих мощностей, освоении международных стандартов, укрепляя
позиции крупного ресурсного хаба, регулирующего региональные цепочки
добавленной стоимости.

Ограниченность официальной статистики по запасам не позволила полностью
количественно оценить потенциал региона {[}16,17{]}, однако
сравнительный анализ ресурсной базы центральноазиатских стран и их
нынешнего геополитического положения подтверждает гипотезу о
затрудненной интеграции региона в ГЦСС из-за геополитической
конкуренции, институциональных рамок и трудностями с переработкой РМЗ.

Незначительный вклад стран Центральной Азии в глобальные и региональные
цепочки создания стоимости объясняется сложностями экономической
диверсификации, высокими затратами на деловые транзакции и
высокорискованной инвестиционной средой, значительно снижающей
коммерческие способности. Низкий промышленный потенциал вынуждает
полагаться на Россию или Китай, что усугубляет инфраструктурные и
инвестиционные риски, что подтверждают результаты PESTEL-анализа,
оказывающих влияние на формирование редкоземельной отрасли (таблица 3).
\end{multicols}

\fig[0.9\textwidth]{e/image4}[Рис.3 - Экспорт и импорт РМЗ Казахстана в 2020-2024 гг.\\\normalfont{\emph{Источник: Бюро национальной статистики Агентства по
стратегическому планированию и реформам Республики Казахстан
{[}Электронный ресурс{]}. - Режим доступа:
\href{https://stat.gov.kz/ru/industries/business-statistics/stat-industrial-production/}{https://stat.gov.kz}}}]

\begin{multicols}{2}
Результаты формализованного PESTEL-анализа позволили ранжировать факторы
влияния на отрасль РМЗ в Центральной Азии (таблица 4). Установлено,
что~политические~и~экологические факторы оказывают критическое
сдерживающее воздействие. Экологический фактор и стандарты ESG выделены
как ключевое «окно возможностей»: внедрение экологически ответственной
добычи является обязательным условием для включения региональных
компаний в цепочки поставок системно значимых экономик (ЕС, США).
Высокая зависимость от технологий КНР и РФ в вопросах разделения
редкоземельных элементов блокирует выход региона на рынки высоких
переделов.
\end{multicols}

\tcap{Таблица 3 - PESTEL-анализ факторов, влияющих на отрасль РМЗ}
\begin{longtblr}[
  label = none,
  entry = none,
]{
  width = \linewidth,
  colspec = {Q[473]Q[469]},
  cells = {font = \footnotesize},
  hlines,
  vlines,
}
{\textbf{P (политические)}\\\textbf{-} рост конкуренции мировых держав за контроль над запасами РЗМ стимулирует интерес к Центральной Азии;\\\textbf{- }ограничения диверсификации\\из-за зависимости Центральной Азии от Китая и России;\\- расширение открытого доступа инвесторов к геологическим данным;\\- установление прозрачных контрактов на добычу РМЗ и лицензий на разведку полезных ископаемых с целью стимулирования предприятий участвовать в ГЦСС;\\- участие в ГЦСС с упором на переработку полезных ископаемых усиливает стратегические региональные позиции;\\- распространение на Центральную Азию конкуренции между Саудовской Аравией и ОАЭ за захват доли рынка переработки РМЗ;\\- создание странами Центральной Азии региональных альянсов.} & {\textbf{E (экономические)}\\\textbf{-} высокая потребность в РЗМ способствует инвестиционной активности, но низкая добавленная стоимость и волатильность цен ограничивают доходы;\\- рост зависимости от зарубежных инвестиций из -за отсутствия собственной инфраструктуры;\\- ограниченность местного предоставления промежуточных ресурсов и услуг вторым и третьим уровнями добавленной стоимости;\\- в среднесрочной и долгосрочной перспективе торговые потоки критически важных материалов не будут подвергаться геополитическому влиянию из-за масштабности их запасов, географической распространенности, возможности переработки во многих местах;\\- сложности получения компаниями доступа к ведущим мировым площадкам для привлечения капитала в горнодобывающую промышленность;\\- привлечение частных юниорских компаний, в том числе иностранных игроков, финансируемых мировыми горнодобывающими компаниями.} \\
{\textbf{S (социальные)}\\- несмотря на достигнутый уровень осведомленности и внедрение стандартов, при высоком спросе на критически важные минералы, их предложение неэластично, а добыча может стать катализатором отрицательных последствий и рисков для местных сообществ;\\- политический режим региона, отсутствие прозрачности в работе правительства, внутрирегиональные споры.} & {\textbf{T (технологические)}\\- отсутствие перерабатывающих мощностей, передовых технологий по разделению и переработки РЗМ, зависимость от китайских и российских технологий;\\- модернизации горнодобывающего сектора посредством развития местных технологических компетенций, цифровых технологий, экологически устойчивых практик;\\- развитие международного научно-исследовательского сотрудничества для содействия в передаче технологий переработки РМЗ, разработки новых карт в прогнозировании залежей стратегических ископаемых и их запасов;\\- \textbf{коллаборации с местными НИИ, }привлечение иностранных компаний к сотрудничеству через ГЧП.} \\
{\textbf{E (Экологические)}\\- возможность отслеживать минеральные продукты по всей цепочке создания стоимости вплоть до их происхождения;\\- освоение стандартов ESG (особенно гидро/солнечная энергия);\\- трудности предприятий в получении экологических разрешений на разведку в регионах присутствия;\\- риски загрязнения (особенно водных ресурсов Кыргызстана, Таджикистана).} & {\textbf{L (Правовые)}\\- совершенствование нормативно-правовой базы, раскрытие информации о производстве и торговли РМЗ для создания равных условий всем субъектам хозяйствования;\\- стремление сохранить конкурентное преимущество в экосистемах в сочетании с растущим спросом на социально- экологические устойчивые договоренности и практики.} 
\end{longtblr}

\tcap{Таблица 4 - Матрица стратегического анализа отрасли РМЗ (PESTEL-результаты)}
\begin{longtblr}[
  label = none,
  entry = none,
]{
  width = \linewidth,
  colspec = {Q[135]Q[631]Q[175]},
  cells = {font = \small},
  row{1} = {c,font=\bfseries\small},
  hlines,
  vlines,
}
Фактор          & Ключевое влияние на ГЦСС                                                           & Степень воздействия \\
Политические    & Рост протекционизма; конкуренция США,ЕС и КНР за доступ к месторождениям           & Критическая         \\
Экономические   & Высокая капиталоемкость добычи и волатильность цен на РЗМ.                         & Высокая             \\
Социальные      & Требования местных сообществ к созданию рабочих мест и прозрачности доходов.       & Средняя             \\
Технологические & Дефицит технологий разделения тяжелых РЗЭ; зависимость от китайского софта         & Высокая             \\
Экологические   & Необходимость сертификации ESG для экспорта на западные рынки                      & Критическая         \\
Правовые        & Несовершенство кодексов о недрах; отсутствие защиты интеллектуальной собственности & Средняя             
\end{longtblr}

\begin{multicols}{2}
В целом, заключаем, что роль центральноазиатских стран в диверсификацию
ГЦСС будет определяться темпами консолидации перерабатывающих мощностей,
финансированием коммерчески жизнеспособных маршрутов. Каждое государство
играет свою роль в ГЦСС, предлагает свое видение экономической
модернизации путем перехода к видам деятельности с более высокой
ценностью, созданию наукоемкого и персонализированного продукта. Переход
от поставок сырья к его глубокой переработке и выпуску продукции более
высоких переделов приведет к увеличению национальной добавленной
стоимости экспорта, что крайне важно для обеспечения долгосрочного
коммерческого успеха перерабатывающих предприятий каждой из стран.

Несмотря на отсутствие технологий для глубокой переработки РЗЭ\emph{,}
инфраструктурные ограничения и геополитические разногласия, страны
Центральной Азии расширяют возможности оптимизации цепочек поставок и
внутрирегиональной переработки.

Для укрепления конкурентных позиций одной лишь диверсификации источников
добычи редкоземельных руд недостаточно. Китайские предприятия по
первичной переработке РЗМ остаются самыми экономически эффективными в
мире благодаря крупномасштабным производствам, потребляющим большие
объемы сырья и поддерживающим рентабельность, более низкой стоимости
рабочей силы, меньшему количеству экологических и нормативных барьеров.

Исследование показало, что индивидуальные попытки стран Центральной Азии
по созданию вертикально-интегрированных производств полного цикла (от
руды до неодимовых магнитов) сталкиваются с проблемой нерентабельности
из-за высоких капитальных затрат. Обоснована экономическая эффективность
модели «региональной кооперации», предполагающей разделение стадий ГЦСС:
концентрация первичной добычи в ресурсных узлах (Кыргызстан,
Таджикистан) с последующим рафинированием на базе технологических хабов
Казахстана и Узбекистана. Такая конфигурация позволяет достичь эффекта
масштаба и снизить транзитные риски в рамках ТМТМ.

Сотрудничество государств Центральной Азии с ЕС, обладающим
технологиями, компетенциями и инвестиционным потенциалом, расширит
возможности как добывающей, так и перерабатывающей отраслей, будет
способствовать развитию местной промышленности путем привлечения
инвестиций в регион, росту социальной и экологической устойчивости,
ускорит доступ государств Центральной Азии к инфраструктуре мирового
класса, исследованиям и инновациям, стимулирующим промышленные
инвестиции {[}18,19{]}.

Для роста устойчивости центральноазиатских стран в перестраивающихся
ГЦСС необходимо:

- укреплять внутренний бизнес-сектор и улучшать деловой климат,
развивать независимую оценку запасов как фундамент доверия инвесторов,
совершенствовать управление ресурсами в соответствии с международными
стандартами ОЭСР, EITI, ESG, TCFD, CRIRSCO/ KAZRC, ICMM, IRMA как основы
прозрачности и доверия на горнорудном рынке; создавать
нормативно-правовые экосистемы, уменьшающие инвестиционные
неопределенности разведки на ранних стадиях;

- углублять региональную интеграцию и трансграничные связи, сохраняя при
этом открытость для глобальной торговли.~~Создавать в регионе надежные
производственные сети и цепочки поставок и стимулировать
внутрирегиональный спрос путем устранения трансграничных ограничений на
перемещение товаров и услуг;

- инвестировать в человеческий капитал и инфраструктуру, ориентированную
на высокотехнологичное производство; защищать права интеллектуальной
собственности и коммерческие интересы; приоритетное внимание к НИОКР для
развития инновационных технологий, повышающих эффективность процессов
разделения и очистки РЗЭ;

- разработать стратегию, сочетающую прочные внутренние позиции
(объединение разведки, переработки, экспорта и промышленной политики,
особенно в части финансирования геологоразведки и создания специфической
инфраструктуры для переработки РМЗ) с углубленной интеграцией и
целенаправленным международным партнерством;

- развивать отраслевые кластеры, стратегические отраслевые центры по
оказанию цифровых и бизнес-услуг.

Учитывая высокую стоимость и сложность рентабельной эксплуатации
перерабатывающих предприятий, определяющее место в управлении цепочками
поставок РЗМ, играет интеграция взаимодополняемости ресурсов. Это
подход, основанный на координации стран с запасами РМЗ и потенциал их
переработки, повысит устойчивость и инвестиционную привлекательность
региона. Благодаря взаимодополняемости и созданию региональных хабов,
страны Центральной Азии смогут поддерживать баланс между Китаем, который
вместе с Индией становится основным фактором роста мировой экономики, и
другими странами.

{\bfseries Выводы.} Проведенное исследование подтверждает гипотезу о том,
что статус Центральной Азии как «пассивного поставщика сырья» может быть
преодолен только через институциональные реформы и технологическую
кооперацию. Основным результатом работы стала систематизация рисков
цепочек поставок и обоснование перехода к проактивной политике
управления ресурсами.

Укрепление стратегических позиций региона в глобальной конкуренции за
РЗМ требует не наращивания объемов добычи, а формирования локальных
технологических компетенций в средней части цепочки создания стоимости
(midstream), что позволит увеличить национальную добавленную стоимость
экспорта на прогнозируемые 15-20\% в долгосрочной перспективе.
Странам Центральной Азии предстоит преодолеть ряд ограничений, от
которых зависят преобразование ресурсных активов в рыночные преимущества
и укрепление позиции в мировом предложении CETM {[}20,21{]}.~

Для реализации комплексного подхода, включающего многовекторную
политику, создание региональных перерабатывающих хабов, партнерство с
ЕС, США и Кореей, цифровизацию и внедрение стандартов ESG/EITI для
привлечения инвестиций, региону предстоит преодолеть выявленные в данном
исследовании риски (политическая зависимость от Китая и России,
экономическая волатильность цен, технологическая отсталость,
экологические угрозы и социальные конфликты).

Предложенные меры реализуемы в среднесрочной перспективе: уже к 2026 г.
Казахстан может активизировать работы на 40 месторождениях, в
Узбекистане к 2028 г. на планируемую мощность выведут 76 проектов,
взаимовыгодные партнерства повысят добавленную стоимость экспорта на
20-30\% и снизят зависимость от Китая.

Своим исследованием авторы привлекают внимание к тому, что странам
Центральной Азии предстоит развивать свой потенциал в области
горнодобывающих технологий, отраслей средней и высшей ступени
переработки и разделения полезных ископаемых посредством стратегического
сотрудничества на коммерческих условиях с зарубежными партнерами.

Будущие исследования могут быть улучшены за счет учета негативного
влияния на коммерческую привлекательность центральноазиатских государств
санкций по отношению к странам, окружающим Центральную Азию. На наш
взгляд, в краткосрочной и среднесрочной перспективе санкции приведут к
снижению экспортных возможностей, высоким транзитным расходам и
высокорискованной инвестиционной среде.
\end{multicols}

\begin{center}
{\bfseries Литература}
\end{center}

\begin{refs}
1. IEA. Global Critical Minerals Outlook.- 2025.- 312 p.
\url{URL:https://www.iea.org/reports/global-critical-minerals-outlook-2025-}
Date of access:18.08 2025.

2. Le H. T. Global economic sanctions and export survival: evidence from
cross-country data // Entrepreneurial Business and Economics
Review.-2022.-Vol.10(1).-P.-22. DOI 10.15678/EBER.2022.100101.

3. Cai M. et al. Towards new Trade in Value-Added (TiVA) indicators that
account for modes of supply in services trade: The TiVA-MoS database //
OECD Trade Policy Papers. -2025.- 35 p. ISSN: 1816-6873.

4. Thakur-Weigold B., Miroudot S. Promoting resilience and preparedness
in supply chains// OECD Trade Policy Papers.- 2024. - 48 p. DOI
10.1787/be692d01-en.

5. Transaction data for evidence-based industrial policy//OECD Science,
Technology and Industry Policy Papers.- 2025.- 42 p. DOI
10.1787/46d7c123-en.

6. Vakulchuk R., Overland I. Central Asia is a missing link in analyses
of critical materials for the global clean energy transition//One
Earth.- 2021. - P.1678--1692.

7. Liu S. L. et al. Global rare earth elements projects: New
developments and supply chains // \emph{Ore} Geology Reviews. - 2023. -
Vol.157: 105428. DOI 10.1016/j.oregeorev.2023.105428.

8. Энергетический переход в Узбекистане URL:
\href{https://www.sfa-oxford.com/lithox/critical-minerals-policy-legislation/all-countries/asia/uzbekistan/\%20-\%20Дата}{https://www.sfa-oxford.com/lithox/critical-minerals-policy-legislation/all-countries/asia/uzbekistan/
- Дата} обращения:19.09 2025

9. Uzbekistan plans to implement rare mineral projects worth \$2.6
billion. URL:
https://invexi.org/press/uzbekistan-plans-to-implement-rare-mineral-projects-worth-2-6-billion/.-
Date of access:12.09.2025).

10. Распоряжение Кабинета Министров Кыргызской Республики от 20 марта
2024 года № 106-р Об утверждении Плана мероприятий по реализации Указа
Президента Кыргызской Республики «О реализации Национального проекта по
добыче полиметаллов и редкоземельных элементов для динамичного развития
экономики Кыргызской Республики» от 22 января 2024 года № 5 (с
изменениями и дополнениями по состоянию на 03.12.2024 г.). URL:
\url{https://cbd.minjust.gov.kg/57-19123/edition/22361/ru}.-Дата
обращения: 03.09.2025.

11. IEA. Global Critical Minerals Outlook.- 2024.
URL:\href{https://iea.blob.core.windows.net/assets/ee01701d-1d5c-4ba8-9df6-abeeac9de99a/GlobalCriticalMineralsOutlook2024.pdf}{https://iea.blob.core.windows.net} .- Date of access: 29.01.2026.

12. World Bank. Mining Sector Review 2024. Washington, DC: World Bank,
2024. -- 45 p. URL: \url{https://www.worldbank.org.-} Date of access:
29.01.2026.

13. The Role of Critical Minerals in Clean Energy Transitions. Paris:
IEA, 2025. -- 287 p. URL:
\href{https://www.iea.org/reports/the-role-of-critical-minerals-in-clean-energy-transitions}{https://www.iea.org} .-
Date of access: 29.01.2026.

14.6Wresearch. Turkmenistan Lithium Market (2025--2031): Industry \&
Analysis. New Delhi: 6Wresearch, 2025. URL:
\url{https://www.6wresearch.com.-} Date of access: 29.01.2026.

15. Global Innovation Index 2025: Innovation at a Crossroads/WIPO.- 2025.
URL:
\url{https://www.wipo.int/web-publications/global-innovation-index-2025/assets/80937/global-innovation-index-2025.-}
Date of access: 13.09.2025.

16. Schwellnus C., Haramboure A., Samek L. Policies to strengthen the
resilience of global value chains: Empirical evidence from the COVID-19
shock // OECD Science, Technology and Industry Policy Papers. -2023. -
43 p. DOI 10.1787/fd82abd4-en.

17. Kazakhstan's Resource Economy: Diversification Through Global Value
Chains / Asian Development Bank and Islamic Development Bank. - 2024. -
120 p. DOI 10.22617/FLS240447.

18. Drost N., Cretti G., Van Giersbergen B. Central Asia emerging from
the shadows: European Union-Central Asia relations in evolving Eurasian
geopolitics // Clingendael. - 2025. - P.1-85. URL:
https://www.clingendael.org/publication/central-asia-emerging-shadows-0.-
Date of access: 13.09.2025.

19. Zakirov B. Role of Uzbekistan in the Rare Earth and Critical
Minerals Economy. - 2025. DOI 10.70735/DEUF7298.

20. Туманшиев Д. А., Исова Л. Т., Жұмақанов Ә. Қазақстанның сирек
кездесетін металдары: сыртқы экономикалық саясат // Л.Н. Гумилев
атындағы Еуразия ұлттық университетінің Хабаршысы. - 2025. - № 2(151). -
б.54-67. DOI 10.32523/2616-6887-2025-151-2-54-67.

21. Saidova F. K., Ilhomov O. Rare Earth Elements and Global Power: The
Economic Geography and Geopolitical Stakes of Strategic Resources //
American Journal of Multidisciplinary Bulletin. - 2025. - Vol.3(4). -
P.189-196. DOI 10.5281/zenodo.15311550.
\end{refs}

\begin{center}
{\bfseries References}
\end{center}

\begin{refs}
1. IEA. Global Critical Minerals Outlook.- 2025.- 312 p.
\url{URL:https://www.iea.org/reports/global-critical-minerals-outlook-2025-}
Date of access:18.08 2025.

2. Le H. T. Global economic sanctions and export survival: evidence from
cross-country data // Entrepreneurial Business and Economics
Review.-2022.-Vol.10(1).-P.-22. DOI 10.15678/EBER.2022.100101.

3. Cai M. et al. Towards new Trade in Value-Added (TiVA) indicators that
account for modes of supply in services trade: The TiVA-MoS database //
OECD Trade Policy Papers. -2025.- 35 p. ISSN: 1816-6873.

4. Thakur-Weigold B., Miroudot S. Promoting resilience and preparedness
in supply chains// OECD Trade Policy Papers.- 2024. - 48 p. DOI
10.1787/be692d01-en.

5. Transaction data for evidence-based industrial policy//OECD Science,
Technology and Industry Policy Papers.- 2025.- 42 p. DOI
10.1787/46d7c123-en.

6. Vakulchuk R., Overland I. Central Asia is a missing link in analyses
of critical materials for the global clean energy transition//One
Earth.- 2021. - P.1678--1692.

7. Liu S. L. et al. Global rare earth elements projects: New
developments and supply chains // \emph{Ore} Geology Reviews. - 2023. -
Vol.157: 105428. DOI 10.1016/j.oregeorev.2023.105428.

8. Jenergeticheskij perehod v Uzbekistane. URL:
https://www.sfa-oxford.com/lithox/critical-minerals-policy-legislation/all-countries/asia/uzbekistan/
- Data obrashhenija:19.09 2025.{[}in Russian{]}

9. Uzbekistan plans to implement rare mineral projects worth \$2.6
billion. URL:
https://invexi.org/press/uzbekistan-plans-to-implement-rare-mineral-projects-worth-2-6-billion/.-
Date of access:12.09.2025).

10. Rasporjazhenie Kabineta Ministrov Kyrgyzskoj Respubliki ot 20 marta
2024 goda № 106-r Ob utverzhdenii Plana meroprijatij po realizacii Ukaza
Prezidenta Kyrgyzskoj Respubliki «O realizacii
Nacional' nogo proekta po dobyche polimetallov i
redkozemel' nyh jelementov dlja dinamichnogo razvitija
jekonomiki Kyrgyzskoj Respubliki» ot 22 janvarja 2024 goda № 5 (s
izmenenijami i dopolnenijami po sostojaniju na 03.12.2024 g.). URL:
\url{https://cbd.minjust.gov.kg/57-19123/edition/22361/ru.-Dataobrashhenija:03.09.2025}
{[}in Russian{]}

11. IEA. Global Critical Minerals Outlook.- 2024.
URL:\href{https://iea.blob.core.windows.net/assets/ee01701d-1d5c-4ba8-9df6-abeeac9de99a/GlobalCriticalMineralsOutlook2024.pdf}{https://iea.blob.core.windows.net}.- Date of access: 29.01.2026.

12. World Bank. Mining Sector Review 2024. Washington, DC: World Bank,
2024. -- 45 p. URL: \url{https://www.worldbank.org.-} Date of access:
29.01.2026.

13. The Role of Critical Minerals in Clean Energy Transitions. Paris:
IEA, 2025. -- 287 p. URL:
\href{https://www.iea.org/reports/the-role-of-critical-minerals-in-clean-energy-transitions}{https://www.iea.org}.-
Date of access: 29.01.2026.

14.6Wresearch. Turkmenistan Lithium Market (2025--2031): Industry \&
Analysis. New Delhi: 6Wresearch, 2025. URL:
\url{https://www.6wresearch.com.-} Date of access: 29.01.2026.

15. Global Innovation Index 2025: Innovation at a Crossroads/WIPO.- 2025.
URL:
\url{https://www.wipo.int/web-publications/global-innovation-index-2025/assets/80937/global-innovation-index-2025} .-
Date of access: 13.09.2025.

16. Schwellnus C., Haramboure A., Samek L. Policies to strengthen the
resilience of global value chains: Empirical evidence from the COVID-19
shock // OECD Science, Technology and Industry Policy Papers. -2023. -
43 p. DOI 10.1787/fd82abd4-en.

17. Kazakhstan's Resource Economy: Diversification Through Global Value
Chains / Asian Development Bank and Islamic Development Bank. - 2024. -
120 p. DOI 10.22617/FLS240447.

18. Drost N., Cretti G., Van Giersbergen B. Central Asia emerging from
the shadows: European Union-Central Asia relations in evolving Eurasian
geopolitics // Clingendael. - 2025. - P.1-85. URL:
https://www.clingendael.org/publication/central-asia-emerging-shadows-0.-
Date of access: 13.09.2025.

19. Zakirov B. Role of Uzbekistan in the Rare Earth and Critical
Minerals Economy. - 2025. DOI 10.70735/DEUF7298.

20. Tumanshiev D. A., Isova L. T., Zhұmaқanov Ә. Қazaқstannyң sirek
kezdesetіn metaldary: syrtқy jekonomikalyқ sajasat // L.N. Gumilev
atyndaғy Eurazija ұlttyқ universitetіnің Habarshysy. - 2025. - № 2(151).
- b.54-67. DOI 10.32523/2616-6887-2025-151-2-54-67.{[}in Kazakh{]}

21. Saidova F. K., Ilhomov O. Rare Earth Elements and Global Power: The
Economic Geography and Geopolitical Stakes of Strategic Resources //
American Journal of Multidisciplinary Bulletin. - 2025. - Vol.3(4). -
P.189-196. DOI 10.5281/zenodo.15311550.
\end{refs}

\begin{info}
\hspace{1em}\emph{{\bfseries Сведения об авторах}}

Аубакирова Г. М.- профессор, Карагандинский технический университет
им.  Абылкаса Сагинова», Караганда, Казахстан, е-mail:
rendykar@gmail.com;

Исатаева Ф.М.- ассоциированный профессор, Карагандинский технический
университет им. Абылкаса Сагинова»,Караганда, Казахстан, е-mail:
isataeva.farida@gmail.com;

Мажитова С.К.- кандидат экономических наук, доцент, Карагандинский
университет Казпотребсоюза,Караганда, Казахстан, е-mail:
Skm19@mail.ru;

Тогайбаева Л.И.- кандидат экономических наук, Карагандинский
технический университет им. Абылкаса Сагинова,Караганда, Казахстан,
е-mail: toglusia@mail.ru;;

Оспанова С.М.- ведущий инженер ТОО "ONYX LS", Астана, Казахстан,
е-mail: ospanovasan@gmail.com;

Дюсекеева А.Т.- доцент Казахский университет технологии и бизнеса
имени К. Кулажанова, кафедры химии, химической технологии и экологии,
Астана, Казахстан, е-mail: dyusekeeva\_at@mail.ru.

\hspace{1em}\emph{{\bfseries Information about authors}}

Aubakirova G. M{\bfseries .-} Professor, Abylkas Saginov Karaganda
Technical University, Karaganda, Kazakhstan, e-mail:
rendykar@gmail.com;

Isataeva F. M.- Associate Professor, Abylkas Saginov Karaganda
Technical University, Karaganda, Kazakhstan,

e-mail: isataeva.farida@gmail.com;

Mazhitova S. K.- Candidate of Economic Sciences, Associate Professor,
Karaganda University of Kazpotrebsoyuz, , Karaganda, Kazakhstan,
e-mail: Skm19@mail.ru;

Togaibaeva L. I.- Candidate of Economic Sciences, Abylkas Saginov
Karaganda Technical University, Karaganda, Kazakhstan, e-mail:
toglusia@mail.ru;

Ospanova S. M.- Lead Engineer, TOO "ONYX LS", Astana, Kazakhstan,
e-mail: ospanovasan@gmail.com;

Dyusekeeva A. T.- Associate Professor K. Kulzhanov Kazakh University
of Technology and Business, Astana, Kazakhstan, e-mail:
dyusekeeva\_at@mail.ru.
\end{info}
